% 第3章3.4.2残量网络与负费用环：三点算例的初始可行流 x=(0,0,3)（顺序 s->a,a->t,s->t），费用12。
% 残量弧（剩余容量，费用）：s->a (3,1)，a->t (2,1)，s->t (1,4)，t->s (3,-4)；
% 环 s->a->t->s 费用 1+1-4=-2，剩余容量 min(3,2,3)=2；推2单位后费用 12-4=8（verify.py）。
\begin{tikzpicture}[>=Stealth,line width=.8pt,font=\fontsize{8}{10}\selectfont,
  v/.style={circle,draw=FigureStroke,fill=FigureBlueFill,line width=.8pt,minimum size=6.5mm,inner sep=0pt,text=FigureInk},
  e/.style={->,draw=FigureStroke,line width=.8pt,shorten >=.5pt},
  lab/.style={text=FigureInk,inner sep=1.5pt},
  note/.style={text=FigureMuted,font=\FigureNoteFont\fontsize{7.5}{9}\selectfont}]
  \begin{scope}
  \node[v] (s) at (0,0) {$s$};
  \node[v] (a) at (1.8,1.4) {$a$};
  \node[v] (t) at (3.6,0) {$t$};
  \draw[e] (s)--node[lab,above left=-1pt]{$0/3$} (a);
  \draw[e] (a)--node[lab,above right=-1pt]{$0/2$} (t);
  \draw[e,draw=FigureBlue,line width=1.3pt] (s)--node[lab,below]{$3/4$} (t);
  \node at (1.8,-1.25) {(a) 可行流：弧旁为$x/u$，费用$12$};
  \end{scope}
  \begin{scope}[xshift=5.4cm]
  \node[v] (s) at (0,0) {$s$};
  \node[v] (a) at (1.8,1.4) {$a$};
  \node[v] (t) at (3.6,0) {$t$};
  \draw[e,draw=FigureGreen,line width=1.3pt] (s)--node[lab,above left=-1pt]{$(3,1)$} (a);
  \draw[e,draw=FigureGreen,line width=1.3pt] (a)--node[lab,above right=-1pt]{$(2,1)$} (t);
  \draw[e] (s) to[bend left=14] node[lab,above,pos=.62]{$(1,4)$} (t);
  \draw[e,draw=FigureGreen,line width=1.3pt] (t) to[bend left=14] node[lab,below]{$(3,-4)$} (s);
  \node[note] at (1.8,2.1) {环$s\to a\to t\to s$的费用为$1+1-4=-2$};
  \node at (1.8,-1.25) {(b) 残量网络：弧旁为（剩余容量，费用）};
  \end{scope}
\end{tikzpicture}
